\documentclass{stp}

\newcommand{\en}[1]{\textit{#1}}

\begin{document}

\stptitle{
  вид        = контрольная,
  номер      = 2,
  тема       = {Интеграция и оптимизация},
  дисциплина = {ИТиВП},
  студент    = {А. С. Образцов},
  группа     = {320601},
  преподаватель = {Н. В. Хаджинова},
}

\section{Цель работы}

Целью контрольной работы является интеграция клиентского приложения на
\en{React} с серверной частью на \en{Node.js}/\en{Express}, оптимизация
производительности клиента и подготовка развёртывания в облаке. Предметная
область --- сервис планирования задач по методу \en{GTD} (\en{Getting Things
Done}) с контекстами и напоминаниями (вариант~27). Серверная часть
(\en{REST API} с \en{JWT}, \en{PostgreSQL}, ролями \en{user}, \en{moderator},
\en{admin}) разработана в контрольной работе~1; в этой работе к ней добавлен
клиент в папке \en{client/}, а основной сущностью интерфейса является задача
(\en{Task}): список, страница задачи, форма создания и редактирования.

\section{Ссылка на приложение}

Развёрнутое приложение: фронтенд --- [будет заполнено после развёртывания],
бэкенд --- [будет заполнено после развёртывания]. Исходный код размещён в
репозитории
\en{https://github.com/\allowbreak{}Alexandr-Obraztsov/\allowbreak{}gtd-task-service},
ветка \en{kr22}. Конфигурации развёртывания (\en{render.yaml},
\en{client/vercel.json}) подготовлены и описаны в
подразделе~\ref{sub:deploy}.

\section{Ход работы}

\subsection{Интеграция клиента и сервера}

Клиент создан на \en{Vite} с библиотеками \en{React}, \en{React Router} и
\en{axios}. Состояние хранится в хуках \en{useState}, \en{useEffect} и
\en{useReducer}, готовых менеджеров состояния и кэша не используется. Запросы
к \en{REST API} собраны в отдельный слой \en{src/api}: единый экземпляр
\en{axios} (\en{http.js}), функции авторизации (\en{auth.js}) и по одной
функции на операцию с задачами (\en{tasks.js}). Страницы не знают адресов и
заголовков: они вызывают \en{listTasks}, \en{getTask}, \en{createTask},
\en{replaceTask} (метод \en{PUT}), \en{patchTask} и \en{deleteTask}.
Базовый адрес берётся из переменной окружения \en{VITE\_API\_URL}, которую
\en{Vite} подставляет при сборке.

Маршруты клиента: главная (\en{/}), вход и регистрация (\en{/login},
\en{/register}), список задач (\en{/tasks}), форма создания
(\en{/tasks/new}), страница задачи (\en{/tasks/:id}) и форма редактирования
(\en{/tasks/:id/edit}). Маршруты задач обёрнуты в компонент
\en{ProtectedRoute}, который отправляет неавторизованного пользователя на
страницу входа.

\subsection{Сервис \en{axios}, токены и обновление сессии}

\en{Access}"=токен хранится только в памяти модуля, а не в
\en{localStorage}, поэтому скрипт, прочитавший хранилище, токен не получит:

\begin{stplisting}
let accessToken = null;

export function getAccessToken() { return accessToken; }
export function setAccessToken(token) { accessToken = token; }
export function clearAccessToken() { accessToken = null; }
\end{stplisting}

\en{Refresh}"=токен сервер кладёт в \en{httpOnly cookie}, клиент его не
видит; чтобы браузер отправлял \en{cookie} на другой источник, экземпляр
\en{axios} создаётся с \en{withCredentials}. Перехватчики добавляют заголовок
\en{Authorization} и при ответе \en{401 TOKEN\_EXPIRED} один раз обновляют
сессию и повторяют запрос:

\begin{stplisting}
export const http = axios.create({
  baseURL: API_URL,
  withCredentials: true,
  timeout: 15000,
});

let refreshPromise = null;

export function refreshSession() {
  if (!refreshPromise) {
    refreshPromise = http
      .post('/auth/refresh', null, { skipAuthRefresh: true })
      .then(({ data }) => {
        setAccessToken(data.accessToken);
        return data;
      })
      .finally(() => { refreshPromise = null; });
  }
  return refreshPromise;
}

http.interceptors.request.use((config) => {
  const token = getAccessToken();
  if (token) config.headers.Authorization = 'Bearer ' + token;
  return config;
});

http.interceptors.response.use(
  (response) => response,
  async (error) => {
    const { config, response } = error;
    const expired = response?.status === 401
      && response.data?.error?.code === 'TOKEN_EXPIRED';
    if (!expired || !config || config.skipAuthRefresh
        || config.retried) {
      return Promise.reject(error);
    }
    try {
      await refreshSession();
    } catch (refreshError) {
      clearAccessToken();
      onSessionExpired();
      return Promise.reject(error);
    }
    return http({ ...config, retried: true });
  },
);
\end{stplisting}

Запрос обновления общий для всех одновременных вызовов
(\en{refreshPromise}). Это принципиально: сервер считает повторное
предъявление отозванного \en{refresh}"=токена кражей и отзывает все сессии, а
\en{React StrictMode} в режиме разработки запускает эффект дважды. После
перезагрузки страницы провайдер аутентификации вызывает тот же
\en{refreshSession}, получает новый \en{access}"=токен по \en{cookie} и
пользователь остаётся в системе:

\begin{stplisting}
useEffect(() => {
  setSessionExpiredHandler(() => dispatch({ type: 'anonymous' }));
  authApi
    .restoreSession()
    .then((user) => dispatch({ type: 'authenticated', user }))
    .catch(() => dispatch({ type: 'anonymous' }));
}, []);
\end{stplisting}

Ошибки \en{axios} приводятся к единому виду \en{\{ status, code, message,
details \}}: ответ сервера, таймаут и отсутствие сети различаются, а ошибки
валидации \en{Joi} (\en{details}) превращаются в подсказки у полей формы.

\subsection{Список задач: загрузка, ошибка, пустой результат}

Страница списка загружает задачи хуком \en{useTasksList}. Он хранит
состояние редьюсером (\en{items}, \en{total}, \en{loading}, \en{error}) и
при смене фильтров отменяет предыдущий запрос через \en{AbortController},
чтобы устаревший ответ не затёр новый:

\begin{stplisting}
useEffect(() => {
  const controller = new AbortController();
  dispatch({ type: 'load' });
  tasksApi
    .listTasks({ q, status, limit, offset }, controller.signal)
    .then((page) => dispatch({ type: 'loaded', page }))
    .catch((error) => {
      if (!isCancelled(error)) {
        dispatch({ type: 'failed', error: normalizeError(error) });
      }
    });
  return () => controller.abort();
}, [q, status, limit, offset, reloadKey]);
\end{stplisting}

Три состояния отображаются явно, в порядке приоритета: ошибка с кнопкой
повтора, индикатор первой загрузки, пустой список (с разным текстом для
пустой базы и для пустого результата поиска) и сам список:

\begin{stplisting}
let body;
if (error) {
  body = <ErrorMessage error={error} onRetry={reload} />;
} else if (loading && items.length === 0) {
  body = <Spinner label="Загружаем задачи…" />;
} else if (items.length === 0) {
  body = filtered
    ? <EmptyState title="Ничего не найдено" hint="…" />
    : <EmptyState title="Задач пока нет" hint="…">
        <Link to="/tasks/new">Создать задачу</Link>
      </EmptyState>;
} else {
  body = (
    <ul>{items.map((t) => <TaskItem key={t.id} … />)}</ul>
  );
}
\end{stplisting}

При повторной загрузке (смена поиска или страницы) прежний список остаётся на
экране с пониженной непрозрачностью, поэтому интерфейс не «мигает» индикатором.
Вид списка представлен на рисунке~\ref{fig:ui-list}, состояние загрузки --- на
рисунке~\ref{fig:ui-loading}, ошибки --- на рисунке~\ref{fig:ui-error},
пустого результата --- на рисунке~\ref{fig:ui-empty}. Загрузка получена
искусственной задержкой ответа списка на 2,5~с, ошибка --- обрывом запроса.

\stpimage[150mm]{img/uiList.png}{Список задач}{fig:ui-list}

\stpimage[150mm]{img/uiLoading.png}{Состояние загрузки списка}{fig:ui-loading}

\stpimage[150mm]{img/uiError.png}{Ошибка загрузки списка с кнопкой повтора}{fig:ui-error}

\stpimage[150mm]{img/uiEmpty.png}{Пустой результат поиска}{fig:ui-empty}

\subsection{Страница задачи и форма}

Страница задачи параллельно запрашивает задачу и список контекстов
(\en{Promise.all}), чтобы показать имя контекста, а ошибка \en{404}
отображается без кнопки повтора. Кнопка удаления видна только роли
\en{admin}, как требует серверная политика. Страница представлена на
рисунке~\ref{fig:ui-details}.

\stpimage[150mm]{img/uiDetails.png}{Страница задачи}{fig:ui-details}

Одна форма служит и для создания, и для редактирования: при наличии
параметра \en{id} она загружает задачу и отправляет \en{PUT} (полную замену),
иначе --- \en{POST}. Перед отправкой выполняется проверка на клиенте, после
отправки ошибки сервера \en{VALIDATION\_ERROR} раскладываются по полям:

\begin{stplisting}
function validate(form) {
  const errors = {};
  if (!form.title.trim()) errors.title = 'Введите название задачи';
  if (form.title.length > 200) {
    errors.title = 'Не более 200 символов';
  }
  if (form.status === 'scheduled' && !form.dueDate) {
    errors.dueDate = 'Для статуса «Запланировано» нужен срок';
  }
  return errors;
}

async function handleSubmit(event) {
  event.preventDefault();
  const found = validate(form);
  setErrors(found);
  if (Object.keys(found).length > 0) return;
  setSaving(true);
  try {
    const saved = isEdit
      ? await replaceTask(id, toPayload(form))
      : await createTask(toPayload(form));
    navigate('/tasks/' + saved.id);
  } catch (error) {
    const normalized = normalizeError(error);
    setErrors(fieldErrors(normalized));
    setSubmitError(normalized);
    setSaving(false);
  }
}
\end{stplisting}

Форма редактирования показана на рисунке~\ref{fig:ui-edit}, ошибка
проверки (статус <<Запланировано>> без срока) --- на
рисунке~\ref{fig:ui-forminvalid}, главная страница с напоминаниями ---
на рисунке~\ref{fig:ui-home}, ответ сервера на неверный пароль --- на
рисунке~\ref{fig:ui-loginerror}. Скриншоты сняты с работающего
приложения (\en{vite dev}, \en{API} и \en{PostgreSQL} на локальной машине)
автоматизированным \en{Chromium}.

\stpimage[150mm]{img/uiEdit.png}{Форма редактирования задачи}{fig:ui-edit}

\stpimage[150mm]{img/uiFormInvalid.png}{Ошибка проверки формы}{fig:ui-forminvalid}

\stpimage[150mm]{img/uiHome.png}{Главная страница с напоминаниями}{fig:ui-home}

\stpimage[150mm]{img/uiLoginError.png}{Вход с неверным паролем}{fig:ui-loginerror}

\subsection{Реализованные оптимизации}

Реализовано четыре разных приёма плюс опрос напоминаний. Для базовых
замеров каждый приём можно отключить параметром адреса (\en{?debounce=0},
\en{?memo=0}, \en{?optimistic=0}); значения читает модуль
\en{perf/flags.js}.

\textbf{Debounce поиска.} Без него каждый введённый символ запускает
запрос к серверу. Хук возвращает значение с задержкой 400~мс, а запрос
зависит от задержанного значения:

\begin{stplisting}
export function useDebouncedValue(value, delay) {
  const [debounced, setDebounced] = useState(value);
  useEffect(() => {
    const timer = setTimeout(() => setDebounced(value), delay);
    return () => clearTimeout(timer);
  }, [value, delay]);
  return debounced;
}

const debouncedSearch = useDebouncedValue(searchInput.trim(), 400);
const q = flags.debounce ? debouncedSearch : searchInput.trim();
\end{stplisting}

Поиск выполняет сервер: в \en{GET /tasks} добавлен параметр \en{q} (подстрока
названия без учёта регистра, спецсимволы \en{LIKE} экранируются).

\textbf{Route-based code splitting.} Каждая страница подключается через
\en{React.lazy}, \en{Suspense} с индикатором лежит в общем макете вокруг
\en{Outlet}, поэтому шапка не пропадает во время загрузки чанка:

\begin{stplisting}
const HomePage = lazy(() => import('./pages/HomePage'));
const TaskListPage = lazy(() => import('./pages/TaskListPage'));
const TaskDetailsPage = lazy(
  () => import('./pages/TaskDetailsPage'));
const TaskFormPage = lazy(() => import('./pages/TaskFormPage'));

<Suspense fallback={<Spinner label="Загружаем страницу…" />}>
  <Outlet />
</Suspense>
\end{stplisting}

\textbf{\en{React.memo}.} Строка списка обёрнута в \en{memo}, а колбэки
родителя (\en{create}, \en{remove}, \en{toggleDone}) стабилизированы
\en{useCallback}. Без этого каждый введённый символ в поле поиска или
быстрого добавления перерисовывает все 20 строк:

\begin{stplisting}
function TaskItemView({ task, contextName, canDelete,
                        onToggle, onRemove }) {
  countRender('TaskItem');
  ...
}

export const TaskItem =
  flags.memo ? memo(TaskItemView) : TaskItemView;
\end{stplisting}

\textbf{Оптимистичное обновление} создания, удаления и отметки <<готово>>
(фрагмент хука \en{useTasksList}). Новая задача появляется в списке сразу
с пометкой <<сохраняется>>, при ошибке сервера изменение откатывается и
показывается сообщение:

\begin{stplisting}
const create = useCallback(async (title) => {
  tempId.current += 1;
  const id = 'tmp-' + tempId.current;
  const draft = { id, title, notes: '', status: 'inbox',
    dueDate: null, contextId: null, pending: true };
  dispatch({ type: 'add', task: draft });
  try {
    const saved = await tasksApi.createTask({ title });
    dispatch({ type: 'replace', id, task: saved });
  } catch (error) {
    dispatch({ type: 'remove', id });
    fail(error);
  }
}, [fail]);

const remove = useCallback(async (id) => {
  const index = itemsRef.current.findIndex((t) => t.id === id);
  const task = itemsRef.current[index];
  dispatch({ type: 'remove', id });
  try {
    await tasksApi.deleteTask(id);
  } catch (error) {
    dispatch({ type: 'restore', index, task });
    fail(error);
  }
}, [fail]);
\end{stplisting}

\textbf{Опрос напоминаний.} На главной странице хук
\en{useDueReminders} раз в 30~с запрашивает \en{GET /tasks/reminders/due},
пропускает опрос в скрытой вкладке и при уходе со страницы очищает таймер
и отменяет запрос.

\subsection{Измерения оптимизаций}

Все измерения выполнены на локальной машине скриптами из папки
\en{reports/kr22/measure} (\en{Playwright}, \en{Chromium}), сырые данные
находятся в \en{reports/kr22/measurements}. Каждый сценарий повторён по
пять (бандл --- семь) раз для каждого варианта, в таблицах указана медиана.

\textbf{Debounce.} Сценарий: в поле поиска вводится слово <<позвонить>>
(девять символов, 100~мс между нажатиями). Число запросов фиксировалось тремя
независимыми способами: событиями \en{Chromium} (\en{Network}), счётчиком
перехватчика \en{axios} и журналом сервера; значения совпали во всех десяти
прогонах (таблица~\ref{tab:debounce}, рисунок~\ref{fig:chart-debounce}).
Без \en{debounce} на сервер уходят запросы с \en{q=п}, \en{q=по}, \ldots,
\en{q=позвонить}; с ним --- один запрос с полным словом. Каждый запрос без
\en{debounce} дополнительно вызывает предварительный запрос \en{CORS}
(\en{OPTIONS}), так как адрес у каждого запроса свой.

\begin{stptable}
  \caption{Запросы \en{GET /tasks?q=} при вводе слова из 9 символов}\label{tab:debounce}
  \begin{tabularx}{\linewidth}{|X|c|c|}\hline
    Показатель & Без \en{debounce} & С \en{debounce} 400~мс\\\hline
    Запросов от браузера & 9 & 1\\\hline
    Запросов в журнале сервера (\en{GET}) & 9 & 1\\\hline
    Запросов \en{OPTIONS} (preflight) & 9 & 1\\\hline
  \end{tabularx}
\end{stptable}

\stpimage[110mm]{img/chartDebounce.png}{Число запросов поиска до и после \en{debounce} (по данным \en{debounce.json})}{fig:chart-debounce}

\textbf{\en{React.memo}.} Измерение выполнено на \en{profiling}"=сборке
(\en{vite build --mode profiling}, \en{react-dom/profiling}), так как в
обычной сборке \en{Profiler} не работает. Число рендеров считает счётчик
в теле компонента, время --- сумма \en{actualDuration} из \en{React Profiler API}
по всем коммитам (таблица~\ref{tab:memo}, рисунок~\ref{fig:chart-memo}).
При вводе в поле быстрого добавления перерисовывается только родитель
(меняется строка ввода), и с \en{memo} ни одна из 20 строк не
перерисовывается. При вводе в поиск остаются 20 рендеров --- это
отображение полученных результатов. Начальная загрузка даёт 40 и 37
рендеров, потому что после получения списка приходят имена контекстов и
строки с контекстом получают новое свойство; на эти рендеры \en{memo} не
влияет. Ненулевое число коммитов с \en{memo} объясняется тем, что сам
родитель и обёртка \en{Profiler} перерисовываются, а строки --- нет.

\begin{stptable}
  \caption{Рендеры строки списка и время коммитов (20 строк)}\label{tab:memo}
  \begin{tabularx}{\linewidth}{|>{\raggedright\arraybackslash}p{48mm}|*{4}{>{\centering\arraybackslash}X|}}\hline
    Сценарий & Рендеров без \en{memo} & Рендеров с \en{memo} & Время без \en{memo}, мс & Время с \en{memo}, мс\\\hline
    Ввод 9 символов в быстрое добавление & 180 & 0 & 27,3 & 0,4\\\hline
    Ввод 9 символов в поиск & 240 & 20 & 33,0 & 2,4\\\hline
    Первая загрузка списка & 40 & 37 & 4,1 & 4,8\\\hline
  \end{tabularx}
\end{stptable}

\stpimage[150mm]{img/chartMemo.png}{Рендеры и время коммитов до и после \en{React.memo} (по данным \en{memo.json})}{fig:chart-memo}

\textbf{Code splitting.} Размер сборки \en{vite build} без разделения
(все страницы импортируются статически) и с \en{React.lazy}
(таблица~\ref{tab:bundle}). Вывод \en{vite} сохранён в
\en{build-baseline.txt} и \en{build-split.txt}. Стартовый код без
разделения --- один файл 341,05~КБ (109,88~КБ после \en{gzip}); с
разделением основной файл 276,87~КБ, общий файл \en{React} 47,08~КБ и
семь файлов страниц размером от 0,29 до 7,21~КБ. Список задач (7,21~КБ)
и форма (3,73~КБ) не загружаются, пока пользователь на них не перейдёт.
Загрузка страницы входа с пустым кэшем: переданных байт \en{JavaScript}
108,6~КБ против 109,2~КБ, первая отрисовка при эмуляции медленной сети
(1,6~Мбит/с, 150~мс, процессор в 4~раза медленнее) 932~мс против 952~мс.
Выигрыш на первом экране небольшой (около 0,5\% байт и 2\% времени): страницы
малы по сравнению с библиотеками (\en{React}, \en{React Router},
\en{axios}), которые нужны сразу. Приём оправдан для роста приложения:
каждая новая страница не увеличивает стартовый объём.

\begin{stptable}
  \caption{Размер \en{JavaScript} без code splitting и с ним}\label{tab:bundle}
  \begin{tabularx}{\linewidth}{|X|c|c|}\hline
    Показатель & Без разделения & С \en{React.lazy}\\\hline
    Файлов \en{JavaScript} в сборке & 1 & 12\\\hline
    Основной файл, КБ (\en{gzip}) & 341,05 (109,88) & 276,87 (89,94)\\\hline
    Страница входа: файлов \en{JS} / КБ по сети & 1 / 109,2 & 4 / 108,6\\\hline
    Первая отрисовка входа, медленная сеть, мс & 952 & 932\\\hline
    Первая отрисовка входа без ограничений, мс & 48 & 48\\\hline
  \end{tabularx}
\end{stptable}

\textbf{Оптимистичное обновление.} Сценарий: под ролью \en{admin} создаётся
задача, отмечается выполненной и удаляется; задержка сети эмулируется через
\en{CDP} (300~мс на каждый запрос). Измерялось время от клика до изменения
в интерфейсе (таблица~\ref{tab:optimistic},
рисунок~\ref{fig:chart-optimistic}). Без оптимизма интерфейс ждёт ответа
(более двух кругов обмена из"=за предварительного запроса \en{CORS}),
с оптимизмом реагирует за один кадр.

\begin{stptable}
  \caption{Время до изменения интерфейса при задержке сети 300~мс, мс}\label{tab:optimistic}
  \begin{tabularx}{\linewidth}{|X|c|c|}\hline
    Операция & Без оптимизма & С оптимизмом\\\hline
    Создание задачи & 830 & 30\\\hline
    Отметка <<готово>> & 815 & 36\\\hline
    Удаление задачи & 831 & 47\\\hline
  \end{tabularx}
\end{stptable}

\stpimage[110mm]{img/chartOptimistic.png}{Время реакции интерфейса до и после оптимистичного обновления (по данным \en{optimistic.json})}{fig:chart-optimistic}

\subsection{Подготовка бэкенда к развёртыванию}\label{sub:deploy-backend}

Сервер уже слушал \en{process.env.PORT}, читал \en{DATABASE\_URL},
\en{JWT\_SECRET}, \en{JWT\_REFRESH\_SECRET} и \en{CORS\_ORIGINS}. При
подключении клиента найдены и исправлены три проблемы.

\begin{stpnum}
  \item Список методов \en{CORS} не содержал \en{PUT}: браузер отклонял
    предварительный запрос на замену задачи. Метод добавлен, проверка
    вынесена в тест \en{test/cors.test.js} (допустимый источник получает
    ответ 204~с \en{PUT}, чужой --- 403).
  \item \en{Refresh}"=\en{cookie} имела \en{SameSite=Strict}: при размещении
    клиента и \en{API} на разных доменах (\en{Vercel} и \en{Render})
    браузер её не отправит. Добавлена переменная \en{COOKIE\_SAMESITE}
    (\en{strict}, \en{lax} или \en{none}); значение \en{none} автоматически
    включает \en{Secure}, в \en{production} \en{Secure} включён всегда.
    Поведение проверяют пять тестов \en{test/refreshCookie.test.js}.
  \item Облачные базы требуют \en{TLS}: добавлена переменная
    \en{DATABASE\_SSL}, включающая \en{ssl} в \en{Sequelize}.
\end{stpnum}

\begin{stplisting}
// src/config/env.js
const cookieSameSite =
  readChoice('COOKIE_SAMESITE',
    ['strict', 'lax', 'none'], 'strict');
...
refreshCookie: Object.freeze({
  sameSite: cookieSameSite,
  secure: cookieSameSite === 'none' || nodeEnv === 'production',
}),
\end{stplisting}

Все семь тестов проходят командой \en{npm test}. Для проверки поиска
добавлен скрипт \en{npm run seed:demo}, создающий 240 демонстрационных
задач.

\subsection{Конфигурации развёртывания}\label{sub:deploy}

Подготовлены файлы для развёртывания <<в один клик>> (бэкенд подготовлен в подразделе~\ref{sub:deploy-backend}); сами развёртывания
на \en{Render} и \en{Vercel} требуют аккаунтов владельца и в рамках этой
работы не выполнялись. Состояние: [будет заполнено после развёртывания].

Файл \en{render.yaml} (\en{Blueprint}) описывает веб"=сервис \en{Node} и
базу \en{PostgreSQL}; строка подключения берётся из базы, секреты
генерирует платформа, а значение \en{CORS\_ORIGINS} задаётся вручную после
получения адреса фронтенда:

\begin{stplisting}
databases:
  - name: gtd-postgres
    plan: free
services:
  - type: web
    name: gtd-task-service
    runtime: node
    buildCommand: npm ci
    startCommand: npm start
    healthCheckPath: /health
    autoDeploy: true
    envVars:
      - { key: NODE_ENV, value: production }
      - { key: TRUST_PROXY, value: "1" }
      - { key: DATABASE_SSL, value: "true" }
      - { key: COOKIE_SAMESITE, value: none }
      - key: DATABASE_URL
        fromDatabase: { name: gtd-postgres,
                        property: connectionString }
      - { key: JWT_SECRET, generateValue: true }
      - { key: JWT_REFRESH_SECRET, generateValue: true }
      - { key: CORS_ORIGINS, sync: false }
\end{stplisting}

Файл \en{client/vercel.json} включает перенаправление всех путей на
\en{index.html} (без него обновление страницы \en{/tasks/5} вернуло бы
\en{404}) и долгое кэширование хешированных файлов:

\begin{stplisting}
{
  "framework": "vite",
  "buildCommand": "npm run build",
  "outputDirectory": "dist",
  "rewrites": [{ "source": "/(.*)", "destination": "/index.html" }],
  "headers": [{ "source": "/assets/(.*)", "headers": [
    { "key": "Cache-Control",
      "value": "public, max-age=31536000, immutable" }] }]
}
\end{stplisting}

Образец переменных клиента находится в \en{client/.env.example}
(\en{VITE\_API\_URL}), порядок развёртывания --- в разделе <<Деплой>>
файла \en{README.md}.

\section{Ответы на контрольные вопросы}

\begin{stpnum}
  \item Интеграция построена на \en{axios} с единым экземпляром
    (\en{baseURL} из \en{VITE\_API\_URL}, \en{withCredentials},
    перехватчики), слоем функций \en{src/api} и хуками \en{useEffect},
    \en{useReducer}. Маршрутизация --- \en{React Router}, сборка ---
    \en{Vite}. Бэкенд --- \en{REST API} на \en{Express}, обмен в формате
    \en{JSON}, ошибки приходят в едином виде с кодом.
  \item Состояния описаны явно в \en{useTasksList} (\en{loading},
    \en{error}, \en{items}) и выводятся компонентами \en{Spinner},
    \en{ErrorMessage} с повтором и \en{EmptyState}; код приведён в
    подразделе <<Список задач>>, экраны --- на
    рисунках~\ref{fig:ui-loading}, \ref{fig:ui-error} и \ref{fig:ui-empty}.
    Ошибки валидации формы показываются у полей.
  \item Реализованы \en{debounce} поиска, \en{route-based code splitting},
    \en{React.memo} и оптимистичные обновления, а также опрос
    напоминаний. Выбор обусловлен узкими местами менеджера задач: поиск
    при вводе порождает лавину запросов, строки списка перерисовываются
    при каждом нажатии клавиши, а создание и удаление должны отзываться
    мгновенно при медленной сети. Виртуальная прокрутка не нужна: страница
    содержит 20 строк.
  \item Тремя способами: счётчиком запросов (события \en{Chromium},
    перехватчик \en{axios}, журнал сервера), \en{React Profiler API}
    на \en{profiling}"=сборке и счётчиком рендеров, выводом \en{vite build}
    и размерами загрузки по \en{CDP} с эмуляцией сети. Результат:
    9 запросов против 1, 180 рендеров против 0, реакция 830~мс против
    30~мс; code splitting дал лишь около 0,5\% байт на первом экране.
    Скрипты и сырые данные приложены к репозиторию.
  \item Фронтенд предполагается разместить на \en{Vercel} с корнем
    \en{client/}: платформа сама определяет \en{Vite}, собирает проект
    командой \en{npm run build} и перевыпускает его при каждом \en{push}
    в подключённую ветку. Конфигурация --- \en{client/vercel.json}.
    Развёртывание: [будет заполнено после развёртывания].
  \item Бэкенд и база описаны в \en{render.yaml}: веб"=сервис
    (\en{npm ci}, \en{npm start}, проверка \en{/health}) и \en{PostgreSQL}.
    Переменные: \en{NODE\_ENV=production}, \en{TRUST\_PROXY=1},
    \en{DATABASE\_URL}, \en{DATABASE\_SSL=true}, \en{JWT\_SECRET},
    \en{JWT\_REFRESH\_SECRET}, \en{COOKIE\_SAMESITE=none},
    \en{CORS\_ORIGINS}; \en{PORT} выдаёт платформа.
    Развёртывание: [будет заполнено после развёртывания].
  \item \en{HTTPS} обеспечивают платформы размещения; \en{cookie} с
    \en{refresh}"=токеном имеет флаги \en{HttpOnly}, \en{Secure} и
    \en{SameSite=None} (разные домены). \en{CORS} разрешает только
    источники из \en{CORS\_ORIGINS} вместе с \en{credentials}, чужие
    получают 403. \en{Access}"=токен живёт 15 минут и хранится только в
    памяти, \en{refresh}"=токен вращается при каждом обновлении, а его
    повтор отзывает все сессии. Весь вывод в \en{React} экранируется
    автоматически, \en{dangerouslySetInnerHTML} не применяется.
  \item \en{Code splitting} --- разделение кода приложения на части,
    которые загружаются по мере необходимости, а не одним файлом.
    \en{React.lazy(() => import('./pages/X'))} превращает динамический
    импорт в компонент, а \en{Suspense} показывает запасной интерфейс,
    пока чанк грузится. Сборщик создаёт для каждой страницы отдельный
    файл (см. таблицу~\ref{tab:bundle}).
  \item Мемоизация сама стоит памяти и времени на сравнение свойств, поэтому
    в большом приложении она может не окупаться. Типичные проблемы:
    новые объекты, массивы и функции в свойствах каждый раз ломают
    \en{memo}; неполный список зависимостей \en{useMemo} и
    \en{useCallback} даёт устаревшие значения; мемоизация усложняет код и
    скрывает настоящие причины лишних рендеров. Применять её нужно после
    измерения: здесь она оправдана замером (180 рендеров против 0).
  \item \en{Debounce} откладывает действие до паузы в событиях и
    выполняет его один раз после последнего (поиск при вводе), а
    \en{throttle} выполняет действие не чаще одного раза за интервал
    (прокрутка, изменение размера окна). Для поиска использован
    \en{debounce}, потому что нужен только итоговый запрос; \en{throttle}
    в приложении не применён.
  \item Виртуальная прокрутка (\en{react-window}) создаёт в \en{DOM}
    только строки, попадающие в окно просмотра, и подменяет их при
    прокрутке; остальные существуют как данные. Она нужна при тысячах
    элементов, когда \en{DOM} становится тяжёлым. В работе не применена:
    список постраничный, по 20 строк.
  \item Библиотеки кэширования (\en{React Query}, \en{SWR}) не
    использовались, ручного кэша ответов нет. Косвенно работают
    \en{Cache-Control: immutable} для файлов сборки на \en{Vercel},
    кэширование ответа \en{preflight} (600~с) и общий запрос обновления
    токена. При развитии стоило бы подключить \en{React Query}.
  \item В \en{Vercel} открывается \en{Settings} $\to$ \en{Environment
    Variables}, добавляется \en{VITE\_API\_URL} со значением адреса
    бэкенда без слеша на конце (например,
    \en{https://gtd-task-service.onrender.com}) для окружений
    \en{Production} и \en{Preview}, после чего проект пересобирается:
    переменные \en{Vite} подставляются при сборке, а не во время работы.
    Доступны только переменные с префиксом \en{VITE\_}, секретов в них
    хранить нельзя. Значения для этой работы: [будет заполнено после
    развёртывания].
  \item Использованы \en{React Profiler API} (компонент \en{Profiler} на
    \en{profiling}"=сборке), счётчик рендеров, \en{Chromium DevTools
    Protocol} для эмуляции сети и процессора, события \en{Network} и
    журнал сервера. Аудит \en{Lighthouse} и расширение \en{React
    DevTools} в работе не запускались; их заменяют перечисленные
    программные измерения.
  \item Оптимистичное обновление применено к созданию, удалению и
    отметке задачи: интерфейс меняется сразу, запрос идёт в фоне, при
    ошибке выполняется откат. Фрагмент кода приведён в подразделе
    <<Реализованные оптимизации>> (функции \en{create} и \en{remove}),
    эффект --- в таблице~\ref{tab:optimistic}.
  \item При большем времени: выполнить развёртывание и аудит
    \en{Lighthouse} на реальных адресах, подключить \en{React Query}
    для кэша и повторных попыток, добавить виртуальную прокрутку и
    бесконечную прокрутку, \en{WebSocket} вместо опроса напоминаний,
    тесты клиента (\en{Vitest}, \en{Testing Library}) и сквозные тесты,
    вынос \en{refresh}"=токена в схему с защитой от \en{CSRF}
    (отдельный заголовок), а также режим офлайн.
\end{stpnum}

\stpunnumbered{Заключение}

В ходе контрольной работы разработан клиент на \en{React} и \en{Vite},
интегрированный с \en{REST API} сервиса планирования задач по методу
\en{GTD}: вход и регистрация с \en{JWT} (токен доступа в памяти, обновление
по \en{httpOnly cookie}), список, страница и форма задачи, явная обработка
загрузки, ошибок и пустого списка.

Реализованы и измерены четыре оптимизации. \en{Debounce} сократил число
запросов поиска с 9 до 1, \en{React.memo} --- число рендеров строк с 180 до 0
при вводе в поле, оптимистичные обновления --- время реакции с 830 до 30~мс
при задержке сети 300~мс. \en{Code splitting} уменьшил основной файл с 341 до
277~КБ, но на первом экране выигрыш составил около 0,5\% байт. Бэкенд
подготовлен к развёртыванию: исправлены \en{CORS} для \en{PUT} и параметры
\en{cookie} для разных доменов, добавлены тесты, \en{render.yaml} и
\en{client/vercel.json}. Развёртывание на облачных платформах и ссылки на
работающее приложение: [будет заполнено после развёртывания].

\end{document}
